\documentclass[11pt,a4paper]{article}
\usepackage[utf8]{inputenc}
\usepackage[T1]{fontenc}
\usepackage[brazil]{babel}
\usepackage{geometry}
\usepackage{hyperref}
\geometry{margin=2.5cm}
\title{EVIDA\\Especificação Científica e Arquitetural\\\large ARCH-001, versão 0.1.0}
\author{Documento de projeto em desenvolvimento}
\date{5 de novembro de 2026}
\begin{document}
\maketitle
\section{Missão}
Construir um sistema científico independente para analisar fontes documentais,
com núcleo C++26, interface web leve e trilha de proveniência reproduzível.
O primeiro corpus é composto por Informativos Conjunturais da Emater/RS-Ascar.
\section{Fronteiras}
A aplicação não representa a instituição que produziu os documentos.
O acervo original permanece imutável; o código e os dados privados são separados.
O sistema não depende do NEXO, MÉTIS ou MORFO para operar.
\section{Epistemologia operacional}
Uma menção documental não equivale a uma ocorrência independente.
Uma coocorrência textual não implica relação causal ou proximidade espacial.
Toda interpretação é separada de sua fonte e da validação humana.
\section{Componentes}
\begin{itemize}
\item \texttt{evida-core}: regras de evidência, proveniência e armazenamento.
\item \texttt{evida-server}: API HTTP local em C++26.
\item \texttt{web/}: interface estática, identidade visual original inspirada no MORFO.
\item \texttt{storage/}: migrações SQL versionadas.
\item \texttt{tests/}: invariantes e testes de integração PROBANTE.
\end{itemize}
\section{Governança da migração}
Preservar o pipeline Python em diretório próprio de referência. Registrar
hash SHA-256 do arquivo histórico. Não importar dados automaticamente. A
migração só será considerada equivalente após comparação de resultados
com a mesma taxonomia, as mesmas fontes e os mesmos parâmetros.
\section{Limites da versão 0.1.0}
Sem ingestão de PDFs, análise de redes, autenticação ou acesso em rede externa.
Escuta apenas em \texttt{127.0.0.1}; banco inicia sem documentos e evidências.
\end{document}
